
This work presents proof-of-concept validation on synthetic data that hierarchical variational autoencoders can embed heterogeneous model collections into unified weight space. Same-task different-architecture models cluster significantly tighter than random baselines (WCSS ratio 0.495, p<0.000001, Cohen's d=1.45) on synthetic model zoo data containing 2,120 models across 2 architecture families and 9 vision tasks. Architecture subspaces exhibit strong compatibility (CKA similarity 0.82 for same-task pairs), and hierarchical pooling preserves 68\% task-relevant information.

These synthetic data results suggest that task-level functional constraints may create architecture-invariant structural features in weight distributions at layer-summary granularity, enabling cross-architecture model property inference without shared base models or hand-designed alignment rules.

Real dataset validation on ModelZooDataset checkpoints is required before publication to confirm external validity and rule out synthetic data artifacts. If real dataset CKA same-task >0.6 AND WCSS Cohen's d >0.5, the hypothesis is validated. If effect size shrinks to d=0.5-0.8 range, hypothesis survives with adjusted claims. If d <0.5 or CKA <0.5, hypothesis mechanism fails on real data.

Future applications pending real validation include model zoo curation (zero-shot task prediction on repositories), cross-architecture transfer learning (task vectors generalizing across CNNs and Transformers), and metadata-free model auditing (training dataset identification from weights alone).

